\section{Selecting polygons near a prescribed maximum}\label{sec:selection}

Put $L=\pi/2$. A \emph{full maximum} is a cap maximizing $P_L$ over
all downward caps. A \emph{joint maximum} is a pair $(U,\alpha)$
maximizing $P_\alpha(U)$ over all caps and all $\alpha\in[L/2,L]$.
We shall construct polygons converging to any prescribed maximum.
This quantifier will also be needed in the equality argument.

\subsection{Compactness and continuity}

\begin{lemma}[Width bounds]\label{lem:selection-width}
If $P_\alpha(U)\ge M/2$ and $L/2\le\alpha\le L$, then
$M\le w_U<6$. After horizontal translation, $U$ lies in
$[-3,3]\times[0,1]$.
\end{lemma}

\begin{proof}
Since $f_U\le1$ and $N_\alpha(U)\ge0$, $P_\alpha(U)\le w_U/2$.
Center the projection and write $w=w_U$. The baseline endpoints give
$\sqrt2 h_U(L/2),\sqrt2 h_U(3L/2)\ge w/2$.
The two walls at $L/2$ therefore give
\[
 N_\alpha(U)(x)\ge(w/2-\sqrt2-|x|)_+.
\]
For $\alpha=L/2$ the terminal first wall dominates this minimum;
for larger $\alpha$ it is one of the visited two-wall minima.
If $w\ge6$, the displayed tent is positive on $[-w/4,w/4]$, and
\[
 P_\alpha(U)\le\frac{(1+\sqrt2)w}{2}-\frac{3w^2}{16}.
\]
This quadratic decreases for $w\ge6$. Its value at six is
$3\sqrt2-15/4<3/4<M/2$, a contradiction. Centering a cap of
width less than six gives the asserted box.
\end{proof}

\begin{lemma}[Continuity of the score]\label{lem:selection-continuity}
On the space of nonempty downward caps in a fixed bounded box,
$P_\alpha(U)$ is jointly continuous in Hausdorff distance and in
$\alpha\in[L/2,L]$.
\end{lemma}

\begin{proof}
Take the box $[-B,B]\times[0,1]$. Its support functions are uniformly
bounded and uniformly Lipschitz in angle. For $x\in[-B,B]$,
\begin{equation}\label{eq:selection-tails}
 (\min(a_t(x),b_t(x)))_+\le2B\min(\tan t,\cot t),\qquad
 (a_\alpha(x))_+\le2B\cot\alpha.
\end{equation}
For example $h_U(t+L)\le B\sin t+\cos t$, so
$b_t(x)\le(B+x)\tan t+1-\sec t\le2B\tan t$.
The first wall gives the other tail bound and the terminal bound.

If two cap supports differ by at most $\eta$, their walls differ by
at most $\eta/\sin\delta$ on $[\delta,L-\delta]$.
The omitted positive tails are at most $2B\tan\delta$. Taking
maxima and minima preserves these estimates. Thus the full niche
changes by at most
\[
 4B\tan\delta+\frac{\eta}{\sin\delta}.
\]
For a varying endpoint below $L$, first restrict to a closed angular
interval away from the axes. Uniform continuity there controls both
the support change and the change of its upper limit. At $\alpha=L$
the second estimate in \eqref{eq:selection-tails} makes the positive
terminal wall tend uniformly to zero. Hence the complete barriers
converge uniformly on the horizontal box.

Projection endpoints converge because they are horizontal supports.
The roofs converge in $L^1$: Hausdorff convergence of bounded convex
sets gives convergence of their indicators almost everywhere away
from the limiting boundary, and integrating the vertical fibers gives
the roof assertion. If the limit is degenerate, the approximating
convex-set areas tend to zero and the same conclusion follows.
The symmetric difference of the moving middle intervals has length
tending to zero. Bounded roofs and barriers now give convergence of
both integrals defining the score.
\end{proof}

\begin{proposition}[Attainment]\label{lem:selection-maximum}
Full and joint maxima exist. A joint maximum can be chosen with its
angle largest among all joint maximizing angles.
\end{proposition}

\begin{proof}
The reference cap has $P_L(\KR)=M/2$. By
\cref{lem:selection-width}, every competitor at least this good can
be centered in $[-3,3]\times[0,1]$. Blaschke's compactness theorem
applies to the nonempty compact convex subsets of this box. Their
downward closure is preserved under Hausdorff limits: approximate
$(x,y)$ by $(x_n,y_n)$ and lower its ordinate to
$\min(y',y_n)$ for $0\le y'\le y$.
The score continuity lemma therefore gives a maximum in the box,
both with $\alpha=L$ fixed and with $\alpha\in[L/2,L]$ variable.
Width bounds remove the box restriction from the conclusion.
The joint maximizing level set is compact; maximize its angle
coordinate to obtain the final assertion.
\end{proof}

\subsection{Canonicalization and used supports}

\begin{lemma}[Chord and height reduction]\label{lem:selection-canonical}
For any positive-width cap $U$ and fixed $0<\alpha\le L$, there
is a cap $U^\sharp$ with the same projection, height one, and an
affine roof on $J_U$, such that $P_\alpha(U^\sharp)\ge P_\alpha(U)$.
Applied to a full or joint maximum, the construction preserves its
score and its angle.
\end{lemma}

\begin{proof}
Let $q$ join the roof values at the endpoints $j_-,j_+$ of $J_U$.
Concavity gives $f_U\ge q$ on $J_U$ and $f_U\le q$ outside it.
For the latter assertion, write $j_-$ as a convex combination of
$x<j_-$ and $j_+$ and solve the concavity inequality for $f_U(x)$;
the right side follows by exchanging endpoints. In particular
$q\ge0$ throughout the projection.
The cap $U^{\rm c}=U\cap\{y\le q(x)\}$ therefore has the same
projection and exterior roof. Inclusion decreases every support,
hence every visited wall and the terminal first wall. Its barrier
decreases, so its score increases.

If $H$ is the height of $U^{\rm c}$, put $\varepsilon=1-H$ and
$U^\sharp=U^{\rm c}+[0,\varepsilon]e_y$.
Its roof rises by $\varepsilon$. Supports in the two visited
families rise by $\varepsilon\sin t$ and $\varepsilon\cos t$;
both raw walls, including the terminal wall, rise by exactly
$\varepsilon$. The maximum with the floor gives
\[
 0\le N_\alpha(U^\sharp)-N_\alpha(U^{\rm c})\le\varepsilon.
\]
The exterior and middle intervals have equal length, so the exterior
gain pays this middle penalty. The projection stays fixed, the
middle roof stays affine, and the height becomes one. Maximality
forces equality of scores when the starting cap is a maximum.
\end{proof}

We call a height-one cap with affine middle roof \emph{canonical}.
Write $f_U(x)=mx+c$ on its middle interval and let $\theta_c$ be
the normal angle of this facet, so $\cot\theta_c=-m$.

\begin{lemma}[Top localization]\label{lem:selection-top}
At every positive full or joint maximum the top face meets $J_U$.
For a canonical maximum the higher middle endpoint has height one.
\end{lemma}

\begin{proof}
Suppose the top face starts at $a>j_+$. Adjoin $(j_+,H)$ and its
downward segment, and take the convex hull. A changed upward support
must have negative horizontal component, because the old top point
$(a,H)$ dominates the new point in all other directions. Its new
support is $\ip{(j_+,H)}{n}$; for $x\in J_U$ and $y\ge0$,
\[
 \ip{(j_+,H)-(x,y)}{n}-1
 =(j_+-x)n_x+(H-y)n_y-1\le0.
\]
Thus changed second walls contribute no positive niche on $J_U$;
first walls and the terminal wall are unchanged. The barrier is
unchanged, whereas the exterior roof strictly increases on
$(j_+,a)$. If the top lies left of $J_U$, insert $(j_-,H)$.
Changed supports then have positive horizontal component and the
same inequality makes their first walls, including the terminal
wall, nonpositive on $J_U$. Both cases contradict maximality.
\end{proof}

\begin{lemma}[Saturation]\label{lem:selection-saturation}
For a canonical joint maximum $(U,\alpha)$,
\begin{equation}\label{eq:selection-saturation}
 \begin{gathered}
 U=U^{\rm sat}\cap\{y\le mx+c\},\\
 U^{\rm sat}=(I_U\times[0,1])\cap
 \bigcap_{0\le t\le\alpha} H_t,\\
 H_t=\{z:\ip z{u(t)}\le h_U(t),\
             \ip z{u^\perp(t)}\le h_U(t+L)\}.
 \end{gathered}
\end{equation}
On either unused open normal arc $(\alpha,L)$ or $(L+\alpha,\pi)$,
the upper curvature measure is zero except possibly at $\theta_c$.
Thus the support is harmonic on the unused arcs away from that normal.
\end{lemma}

\begin{proof}
The saturated cap contains $U$, retains its baseline, and has the
same support in each used direction: its defining inequalities give
one bound and inclusion gives the opposite bound. In particular it
has the same height and complete barrier. Maximality makes its
exterior roof equal to $f_U$, since a strict increase at an interior
exterior point would increase an area integral. Continuity gives
equality at $j_\pm$; cutting to their chord now proves
\eqref{eq:selection-saturation}.
Every boundary point of $U^{\rm sat}$ has an active constraint from
the compact family of defining normals. At arclength almost every
boundary point its outer normal is unique, so equals an active
normal. There is consequently no surface-area measure on unused
open arcs. The chord cut only adds its facet normal $\theta_c$.
The distributional identity $h''+h=0$ on the remaining arcs is the
stated harmonicity.
\end{proof}

\subsection{A penalty that retains any chosen maximum}

Fix a full maximum $U_*$, or a joint maximum $(U_*,\alpha)$ and
keep this angle fixed. Choose a horizontal box with unit slack
beyond both projection endpoints. For $n=2^k$ set
$t_j=j\alpha/n$ and use the sampled barrier
\[
 N_{k,\alpha}(V)=\max\{0,
       \max_{0<j<n}\min(a_{t_j},b_{t_j}),a_\alpha\}.
\]
For $\alpha=L$ the last wall is redundant. Let $P_{k,\alpha}$ be
the corresponding score. Sample outer normals on the two source
grids and on an independent dyadic grid of the whole upper semicircle:
\[
 \Gamma_k=\{jL/n:0\le j\le2n\}
       \cup\{t_j,L+t_j:0\le j\le n\}.
\]
These sets are nested, contain the axes and top, and have dense union.

\begin{lemma}[Uniform grid error and circumscription]\label{lem:selection-grid}
On the fixed boxed cap domain,
\[
 0\le P_{k,\alpha}(V)-P_\alpha(V)\le e_k,
 \qquad e_k\longrightarrow0.
\]
Circumscription by the normals $\Gamma_k$ preserves every sampled
support, projection, and sampled barrier, and increases the exterior
roof. It therefore increases $P_{k,\alpha}$.
\end{lemma}

\begin{proof}
On $[\delta,L-\delta]$ the wall functions have a common modulus of
continuity in angle. A mesh of size $\alpha/2^k$ approximates their
maximum uniformly there. Equation \eqref{eq:selection-tails} makes
the omitted endpoint contributions uniformly small. The terminal
wall is already included exactly. First choose $\delta$ and then
the mesh; this gives uniform barrier approximation. Multiplying
its error by the bounded middle length proves the score estimate.
The circumscribed intersection contains the original cap and retains
each sampled support: an old attaining point remains in it, and the
corresponding defining halfplane gives the upper bound. The axis
supports preserve the projection, and all barrier supports are sampled.
\end{proof}

For $s\in\Gamma_m$ give its level-$m$ sample weight
$\nu_m=2^{-m-1}/|\Gamma_m|$, and define the persistent penalty
\begin{equation}\label{eq:selection-penalty}
 D_k(V,U_*)=\sum_{m=0}^k\nu_m\sum_{s\in\Gamma_m}
                   (h_V(s)-h_{U_*}(s))^2.
\end{equation}
The total weight is at most one. A normal first sampled at an early
level keeps that positive weight at every later level. In particular
\begin{equation}\label{eq:selection-persistence}
 \nu_m(h_V(s)-h_{U_*}(s))^2\le D_k(V,U_*)
 \quad(s\in\Gamma_m,\ m\le k).
\end{equation}

\begin{theorem}[Selection at a prescribed maximum]\label{thm:selection-prescribed}
There are levels $k_i\to\infty$, coefficients $\eta_i\downarrow0$,
and grid polygons $U_i$ converging in Hausdorff distance to the
specified cap $U_*$, such that $U_i$ maximizes
\[
 P_{k_i,\alpha}(V)-\eta_iD_{k_i}(V,U_*)
\]
over every cap in the slack box. No middle affinity or regularity
is required of the selected polygons.
\end{theorem}

\begin{proof}
Choose $\eta_i=1/(i+1)$ and $k_i\ge i$ so that $e_{k_i}\le\eta_i^2$.
Compactness and continuity give a boxed maximum of the penalized
objective. Circumscription increases its score and preserves its
penalty, so a maximizing cap can be taken to be a grid polygon.
The circumscription of $U_*$ has zero penalty and sampled score at
least $P_\alpha(U_*)$. Hence
\[
 P_\alpha(U_*)\le P_{k_i,\alpha}(U_i)-\eta_iD_{k_i}(U_i,U_*)
 \le P_\alpha(U_*)+e_{k_i}-\eta_iD_{k_i}(U_i,U_*),
\]
where the second inequality uses the maximality of $U_*$ at its
fixed angle. Thus $D_{k_i}(U_i,U_*)\le\eta_i\to0$.
Any Hausdorff cluster cap has, by \eqref{eq:selection-persistence},
the same support as $U_*$ at every sampled normal. Density gives
equality on the upper semicircle. For a downward cap, downward
supports are determined by its baseline and horizontal endpoints;
the upper supports consequently determine the entire cap. Every
cluster cap is $U_*$, which proves convergence of the whole sequence.
Its artificial box sides are inactive for all sufficiently large $i$.
\end{proof}

If all support values have absolute value at most $B$, changing
sampled supports by at most $r$ changes the penalty by at most
$4Br+r^2$. If only one normal changes, multiply this bound by its
total persistent weight. These bounds tend to zero after multiplication
by $\eta_i$ and will account for the error in the polygon variations.
